\documentclass[11pt]{article}

\usepackage[letterpaper,margin=1in]{geometry}
\usepackage{times}
\usepackage{graphicx}
\usepackage{booktabs}
\usepackage{multirow}
\usepackage{array}
\usepackage{xcolor}
\usepackage{amsmath}
\usepackage{amssymb}
\usepackage{siunitx}
\usepackage{natbib}
\usepackage[hidelinks]{hyperref}

\sisetup{detect-all}

\title{Detection-Guided Semantic 3D Phenotyping of Cotton Bolls from UAV Imagery Before and After Defoliation}

\author{
Harshitha Manjunatha$^{1}$ \\
$^{1}$Affiliation to be added \\
\texttt{email@domain.edu}
}

\date{}

\begin{document}
\maketitle

\begin{abstract}
Cotton boll phenotyping from unmanned aerial vehicle (UAV) imagery is usually treated as a two-dimensional detection or counting problem, although agronomic decisions often depend on organ-level morphology and visibility. This paper studies whether detection-guided multi-view reconstruction can move cotton boll analysis from image-level counts to three-dimensional morphology. The proposed pipeline combines weak boll detections, mask refinement, camera-pose estimation, dense semantic features, and multi-view association to estimate boll count, three-dimensional location, diameter, volume, visibility, and occlusion before and after defoliation. The pre- and post-defoliation flights provide a paired visibility intervention: the same field was imaged under dense-canopy and exposed-boll conditions, enabling a direct evaluation of how defoliation changes boll recoverability and measurement stability. The manuscript will report reconstruction quality, correspondence quality, boll recovery, morphology accuracy, and robustness across reconstruction, segmentation, and association settings.
\end{abstract}

\textbf{Keywords:} precision agriculture; cotton phenotyping; UAV imagery; 3D reconstruction; cotton boll morphology; semantic correspondence; defoliation; DINOv2; SAM

\section{Introduction}

Cotton yield and harvest readiness are strongly tied to boll development, boll opening, and canopy visibility. UAV-based sensing has made field-scale cotton monitoring practical, but most operational pipelines still reduce the crop to two-dimensional counts, canopy indices, or coarse plant-height summaries. These measurements are useful, yet they do not directly answer organ-level questions: where are the bolls in three-dimensional space, how large are they, how completely are they visible, and how much does defoliation improve their recoverability?

This work treats defoliation as a controlled visibility intervention rather than a nuisance condition. Pre-defoliation imagery captures cotton bolls under dense foliage and severe occlusion, while post-defoliation imagery exposes a larger fraction of boll structure. The paired flights make it possible to quantify how visibility affects detection, reconstruction, and morphology estimation.

The central challenge is that open cotton bolls are difficult for classical photogrammetry. Cotton lint is bright, weakly textured, repetitive, and often partially occluded. As a result, photometric feature matching can be unstable exactly where boll-level reconstruction is needed most. The proposed system therefore uses detections from prior cotton boll counting work as object-level priors, and combines them with mask refinement, camera geometry, and semantic feature correspondences.

\paragraph{Contributions.}
This paper will make four focused contributions:
\begin{enumerate}
    \item A detection-guided 3D phenotyping architecture for cotton bolls from UAV imagery, designed to avoid dense manual annotation.
    \item A pre/post-defoliation evaluation protocol that measures boll visibility, recoverability, and morphology stability.
    \item A multi-view association and localization pipeline that links two-dimensional boll candidates across overlapping views.
    \item A robustness-oriented evaluation covering correspondence quality, reconstruction quality, boll recovery, diameter, volume, and visibility.
\end{enumerate}

\paragraph{Scope.}
The core contribution is computer vision and geometry. A language-model report generator may be used after morphology extraction to summarize quantitative outputs, but it is not part of the reconstruction method and is not required for the paper's main claim.

\section{Related Work}

\subsection{Cotton Boll Detection and Yield Estimation}

Prior cotton monitoring studies have used UAV imagery and image-processing or deep-learning models to detect open bolls and estimate yield. The accepted prior work from this project established a low-annotation counting pipeline for pre- and post-defoliation imagery. In the present paper, that detector is reused as weak supervision: detections provide candidate boll centers and boxes, while the new contribution concerns multi-view association, 3D localization, and morphology estimation.

\subsection{3D Plant Phenotyping}

Three-dimensional plant phenotyping has been studied through structure-from-motion, multi-view stereo, LiDAR, neural radiance fields, and Gaussian splatting. Cotton-specific work has already shown that UAV images can support in-situ boll reconstruction and that 3D Gaussian splatting can map cotton boll and plant architecture. These papers make it important to position the present work carefully: the novelty is not simply that cotton can be reconstructed in 3D, but that detection-guided semantic correspondences are evaluated for boll-level morphology under a paired defoliation intervention.

\subsection{Foundation Features for Correspondence}

Classical reconstruction systems rely on local texture and photometric consistency. Cotton lint creates a difficult case because the target organ is bright and repetitive. Dense self-supervised features such as DINOv2 provide patch-level semantic representations that may remain stable when local texture is weak. This paper evaluates such features as support for multi-view boll association and reconstruction, alongside classical baselines.

\subsection{Promptable Segmentation}

Promptable segmentation models such as SAM and SAM2 can refine object masks from sparse prompts, points, or boxes. Here, promptable segmentation is treated as a mask-refinement stage seeded by the existing boll detector. This keeps the system scalable: only a small validation subset needs manual inspection, rather than dense annotation of thousands of bolls per image.

\section{Method}

\subsection{Problem Definition}

Let $\mathcal{I}^{pre}$ and $\mathcal{I}^{post}$ denote UAV image sets captured before and after defoliation. The goal is to estimate a set of boll instances
\begin{equation}
\mathcal{B} = \{(\mathbf{x}_j, d_j, V_j, \nu_j, o_j)\}_{j=1}^{N},
\end{equation}
where $\mathbf{x}_j \in \mathbb{R}^3$ is the 3D boll center, $d_j$ is diameter, $V_j$ is volume, $\nu_j$ is visibility, and $o_j$ is occlusion.

\subsection{System Overview}

The architecture has eight stages: dataset audit, 2D boll detection, mask refinement, camera-pose estimation, semantic feature extraction, multi-view association, 3D localization, and morphology extraction. The detector supplies weak object priors; geometry supplies scale and pose; semantic features improve association in texture-poor regions.

\subsection{Detection-Guided Candidate Generation}

The existing cotton boll detector is applied to every image. For each image $I_i$, it returns candidate regions
\begin{equation}
\mathcal{D}_i = \{(c_{ik}, b_{ik}, a_{ik}, p_i)\}_{k=1}^{K_i},
\end{equation}
where $c_{ik}$ is the candidate center, $b_{ik}$ is the bounding box, $a_{ik}$ is the detected area, and $p_i$ is the inferred phase. These detections are not assumed to be perfect; they are treated as candidate proposals for later geometric filtering.

\subsection{Mask Refinement}

Each candidate box or center is refined into an approximate boll mask. Two modes will be evaluated: the classical morphology mask from the counting pipeline and a promptable SAM/SAM2 mask initialized from detector boxes. The output mask $M_{ik}$ provides a silhouette used for size estimation and multi-view consistency checks.

\subsection{Camera Geometry and Semantic Correspondence}

Camera poses are estimated with a COLMAP/SfM baseline. In parallel, dense DINOv2 features are extracted from each image. For overlapping image pairs, patch correspondences are obtained by mutual nearest-neighbor search in normalized embedding space:
\begin{equation}
s(u,v) =
\frac{f_i(u)^\top f_l(v)}
{\lVert f_i(u) \rVert_2 \lVert f_l(v) \rVert_2}.
\end{equation}
Matches are filtered by similarity thresholding and geometric verification.

\subsection{Multi-View Boll Association}

Candidate bolls are linked across views using a score that combines geometric consistency, semantic similarity, and mask agreement:
\begin{equation}
A(k,l) =
\lambda_g A_g(k,l) +
\lambda_s A_s(k,l) +
\lambda_m A_m(k,l),
\end{equation}
where $A_g$ measures epipolar or reprojection consistency, $A_s$ measures DINOv2 feature similarity inside the mask, and $A_m$ measures projected mask agreement. Associations below threshold are rejected.

\subsection{3D Localization and Morphology}

High-confidence associated detections are triangulated to estimate boll centers. Diameter is estimated from the multi-view mask extent after converting image scale to metric scale through camera geometry. Volume is reported using two levels of rigor: an ellipsoid approximation for all high-confidence instances, and a visual-hull or point-cluster estimate where sufficient multi-view evidence exists:
\begin{equation}
V_j = \frac{4}{3}\pi a_j b_j c_j.
\end{equation}
Visibility is computed as the fraction of expected overlapping views in which the boll is detected with sufficient mask support.

\section{Experiments}

\subsection{Dataset}

The dataset contains 1,549 valid UAV RGB images captured on 29 September 2025. The configured audit reports 680 pre-defoliation images and 869 post-defoliation images after excluding macOS resource-fork files. Post-defoliation folders contain duplicate basenames that must be handled during reconstruction and evaluation.

\subsection{Evaluation Questions}

The experiments are organized around five questions:
\begin{enumerate}
    \item Does semantic correspondence improve matching in texture-poor cotton regions?
    \item Does post-defoliation improve 3D boll recoverability?
    \item How sensitive are diameter and volume estimates to reconstruction and mask quality?
    \item Which components are necessary for robust multi-view boll association?
    \item Can the method scale from a clean subset to larger field sequences without manual annotation?
\end{enumerate}

\subsection{Main Reconstruction Comparison}

\begin{table}[t]
\centering
\caption{Main 3D reconstruction and boll recovery comparison. Best values will be bolded after experiments are complete. BRR denotes boll recovery rate. RC denotes reconstruction completeness.}
\label{tab:main_reconstruction}
\begin{tabular}{llrrrrr}
\toprule
Category & Method & Registered & Points (K) & RC $\uparrow$ & BRR $\uparrow$ & NN noise $\downarrow$ \\
\midrule
Classical & COLMAP-SIFT Pre & TBD & TBD & TBD & TBD & TBD \\
Classical & COLMAP-SIFT Post & TBD & TBD & TBD & TBD & TBD \\
Learned & DUSt3R/MASt3R Post Subset & TBD & TBD & TBD & TBD & TBD \\
Ours & Detection + COLMAP & TBD & TBD & TBD & TBD & TBD \\
Ours & Detection + COLMAP + DINOv2 & TBD & TBD & TBD & TBD & TBD \\
\bottomrule
\end{tabular}
\end{table}

\subsection{Robustness Grid}

\begin{table}[t]
\centering
\caption{Robustness grid for detection-guided 3D boll phenotyping. The table is designed to report the kind of component-level robustness expected in top-tier vision papers.}
\label{tab:robustness_grid}
\begin{tabular}{llrrrrr}
\toprule
Block & Variant & BRR $\uparrow$ & Diam. MAE $\downarrow$ & Vol. err. $\downarrow$ & Visibility $\uparrow$ & Runtime $\downarrow$ \\
\midrule
Detector & Classical count prior & TBD & TBD & TBD & TBD & TBD \\
Detector & Prior + SAM prompt & TBD & TBD & TBD & TBD & TBD \\
Features & SIFT only & TBD & TBD & TBD & TBD & TBD \\
Features & DINOv2-S/14 & TBD & TBD & TBD & TBD & TBD \\
Features & DINOv2-B/14 & TBD & TBD & TBD & TBD & TBD \\
Features & DINOv2-L/14 & TBD & TBD & TBD & TBD & TBD \\
Association & Geometry only & TBD & TBD & TBD & TBD & TBD \\
Association & Geometry + semantic & TBD & TBD & TBD & TBD & TBD \\
Association & Geometry + semantic + mask & TBD & TBD & TBD & TBD & TBD \\
\bottomrule
\end{tabular}
\end{table}

\subsection{Pre/Post Defoliation Morphology}

\begin{table}[t]
\centering
\caption{Pre- and post-defoliation morphology statistics. Report mean, standard deviation, confidence interval, and effect size for each trait.}
\label{tab:pre_post_morphology}
\begin{tabular}{lrrrrr}
\toprule
Trait & Pre mean $\pm$ SD & Post mean $\pm$ SD & $\Delta$ (\%) & Effect size & $p$ value \\
\midrule
Boll count & TBD & TBD & TBD & TBD & TBD \\
3D recovered bolls & TBD & TBD & TBD & TBD & TBD \\
Diameter (mm) & TBD & TBD & TBD & TBD & TBD \\
Volume (mm$^3$) & TBD & TBD & TBD & TBD & TBD \\
Visibility & TBD & TBD & TBD & TBD & TBD \\
Occlusion & TBD & TBD & TBD & TBD & TBD \\
\bottomrule
\end{tabular}
\end{table}

\subsection{Qualitative Figures}

The final paper should include: (1) pipeline architecture, (2) SIFT versus DINOv2 correspondences on cotton lint, (3) pre/post mask overlays, (4) point-cloud or 3DGS reconstruction comparison, and (5) boll morphology visualization with diameter axes and visibility rays.

\section{Discussion}

The discussion should separate measured evidence from intended use. The expected strength of the method is not perfect reconstruction of every boll in a dense field, but reliable morphology estimation for high-confidence multi-view instances and a quantified understanding of where visibility limits the pipeline. The paper should report failure modes explicitly: repeated white lint texture, wind motion, duplicate frames, low-baseline views, severe leaf occlusion, and scale uncertainty.

The method should also be positioned as a bridge between prior two-dimensional counting work and future field-scale 3D crop intelligence. If the optional reporting layer is included, it should be described as a downstream summary of measured traits, not as a substitute for geometric estimation or agronomic validation.

\section{Conclusion}

This paper proposes a detection-guided semantic 3D phenotyping framework for cotton bolls from UAV imagery. By combining prior 2D boll detection, mask refinement, camera geometry, semantic correspondence, and multi-view association, the system is designed to estimate boll count, 3D location, diameter, volume, visibility, and occlusion before and after defoliation. The final manuscript will quantify whether defoliation improves boll recoverability and whether semantic features provide measurable gains over classical photogrammetric matching in texture-poor cotton regions.


\bibliographystyle{plainnat}
\bibliography{references}

\end{document}
